\section{Vector space and subsapce}
        \subsection{Definition of vector space}
            \hskip \parindent
            \par A \textbf{vector space} $V$ over a field $\mathbb{F}$ consists of the following properties:
            \begin{itemize}
                \item A field $\mathbb{F}$ of scalars.
                \item A set $V$ of objects, called vectors.
                \item An operation, called vector addition $(+: V \times V \rightarrow V)$
                    \begin{itemize}
                        \item [$A_1$]: $(u+v)+w = u+(v+w), \forall u,v,w \in V$
                        \item [$A_2$]: $u+v = v+u, \forall u,v \in V$
                        \item [$A_3$]: $\exists 0 \in V$ s.t. $0+u = u, \forall u \in V$
                        \item [$A_4$]: $\forall u \in V, \exists -u \in V$ s.t. $u + (-u) = 0$
                    \end{itemize}
                \item An operation, called scalar multiplication $(\cdot: \mathbb{F} \times V \rightarrow V)$
                    \begin{itemize}
                        \item [$P_1$]: $1 \cdot v = v, \forall v \in V$
                        \item [$P_2$]: $\alpha(\beta v) = (\alpha \beta)v, \forall \alpha,\beta \in \mathbb{F}, \forall v \in V$
                        \item [$P_3$]: $\alpha(u+v) = \alpha u + \alpha v, \forall \alpha \in \mathbb{F}, \forall u,v \in V$
                        \item [$P_4$]: $(\alpha + \beta)v = \alpha v + \beta v, \forall \alpha,\beta \in \mathbb{F}, \forall v \in V$
                    \end{itemize}
            \end{itemize}
            \par By the way, to prove that a set $V$ is a vector space, it's better to prove the addition and scalar multiplication are closed. 
        \subsection{Examples of vector spaces}
            \subsubsection{$\mathbb{F}^n$ over $\mathbb{F}$}
                \hskip \parindent
                \par $\mathbb{F}^n$ over $\mathbb{F}$
                \begin{equation*}
                    \begin{pmatrix} x_1 \\ x_2 \\ \vdots \\ x_n \end{pmatrix} + \begin{pmatrix} y_1 \\ y_2 \\ \vdots \\ y_n \end{pmatrix} = \begin{pmatrix} x_1 + y_1 \\ x_2 + y_2 \\ \vdots \\ x_n + y_n \end{pmatrix}, \quad \alpha \cdot \begin{pmatrix} x_1 \\ x_2 \\ \vdots \\ x_n \end{pmatrix} = \begin{pmatrix} \alpha x_1 \\ \alpha x_2 \\ \vdots \\ \alpha x_n \end{pmatrix}
                \end{equation*}
            \subsubsection{$V = \{0\}$ over $\mathbb{F}$}
                \hskip \parindent
                \par $V = \{0\}$ over $\mathbb{F}$
                \begin{equation*}
                    0 + 0 = 0, \quad \alpha \cdot 0 = 0, \forall \alpha \in \mathbb{F}
                \end{equation*}
            \subsubsection{$V = \mathbb{F}[x]$ over $\mathbb{F}$}
                \hskip \parindent
                \par $V = \mathbb{F}[x]$ over $\mathbb{F}$
                \begin{equation*}
                    \mathbb{F}[x] = \{p(x) = a_0 + a_1 x + a_2 x^2 + \cdots + a_n x^n \mid a_1, a_2, \ldots, a_n \in \mathbb{F}\}
                \end{equation*}
                \begin{equation*}
                    p(x) = \sum_{i=0}^{n} a_i x^i, \quad q(x) = \sum_{i=0}^{m} b_i x^i
                \end{equation*}
                \begin{equation*}
                    p(x) + q(x) = \sum_{i=0}^{\max \{ n,m \}} (a_i + b_i) x^i
                \end{equation*}
                \begin{equation*}
                    \alpha \cdot p(x) = \sum_{i=0}^{n} \alpha a_i x^i
                \end{equation*}
            \subsubsection{$\mathbb{F}^{m\times n}$}
                \hskip \parindent
                \par The space of $m \times n$ matrices over the field $\mathbb{F}$, denoted $\mathbb{F}^{m\times n}$ (or $M_{m\times n}(\mathbb{F})$).
                \begin{equation*}
                    A, B \in \mathbb{F}^{m \times n}, \quad (A+B)_{ij} = a_{ij} + b_{ij}
                \end{equation*}
                \begin{equation*}
                    \alpha \in \mathbb{F}, \quad (\alpha A)_{ij} = \alpha a_{ij}
                \end{equation*}
            \subsubsection{function f}
                \hskip \parindent
                \par The space of functions from a set $S$ to a field $\mathbb{F}$. $V = \{ f \mid f: S \rightarrow \mathbb{F} $,f is a function$\}$ over $\mathbb{F}$
                \begin{equation*}
                    +: V \times V \rightarrow V, \quad f+g: S \rightarrow \mathbb{F}, \quad (f+g)(x) = f(x) + g(x)
                \end{equation*}
                \begin{equation*}
                    \cdot: \mathbb{F} \times V \rightarrow V, \quad \alpha \cdot f: S \rightarrow \mathbb{F}, \quad (\alpha \cdot f)(x) = \alpha f(x)
                \end{equation*}
        \subsection{Definition of subspace}
            \hskip \parindent
            \par Let $V$ be a vector space over a field $\mathbb{F}$. A subspace of $V$ is a non-empty subset $W$ of $V$ which is itself a vector space over $\mathbb{F}$ with the operations of vector addition and scalar multiplication defined on $V$. We denote it as $W \le V$.
        \subsection{Example of subspace}
            \begin{equation*}
                V = \mathbb{R}^3, \quad W = \{ (x,y,0) \mid x,y \in \mathbb{R} \}, \quad W \le V
            \end{equation*}
        \subsection{Proposition: test for subspace}
            \hskip \parindent
            \par A non-empty subset $W$ of $V$ is a subspace of $V$ if and only if for each pair of vectors $u,v$ in $W$ and each scalar $\alpha$ in $\mathbb{F}$ satisfies:
            \begin{itemize}
                \item [i]: $0 \in W$
                \item [ii]: $u+v \in W$
                \item [iii]: $\alpha v \in W$
            \end{itemize}
        \subsection{Examples of the test for subspace}
            \subsubsection{The solution set of a homogeneous linear system}
                \hskip \parindent
                \par Let $W$ be the solution set of a linear system $Ax=0_m$ with $A \in M_{m \times n}(\mathbb{F})$. Then $W$ is a subspace of $\mathbb{F}^n$.
                \begin{itemize}
                    \item [i]: $0_n \in W: A0_n = 0_m$
                    \item [ii]: $v,w \in W, \text{ then } A(v+w) = Av + Aw = 0_m + 0_m = 0_m, \text{ so } v+w \in W$
                    \item [iii]: $\alpha \in \mathbb{F}, v \in W: \text{ then } A(\alpha v) = \alpha (Av) = \alpha (0_m) = 0_m, \text{ so } \alpha v \in W$
                \end{itemize}
                \par Thus $W$ is a subspace of $\mathbb{F}^n$ ($W \le \mathbb{F}^n$).
            \subsubsection{$A\in \mathbb{F}^{n \times n}$ is symmetric}
                \hskip \parindent
                \par Let $V = \mathbb{F}^{n \times n}$, $W = \{ A \in \mathbb{F}^{n \times n} \mid a_{ij} = a_{ji}, \forall 1 \le i, j \le n \}$.
                \par $W \le V$:
                \begin{itemize}
                    \item [i]: $0_{n \times n} \in W$
                    \item [ii]: $A,B \in W, \text{ then } (A+B)_{ij} = a_{ij} + b_{ij} = a_{ji} + b_{ji} = (A+B)_{ji}, \forall 1 \le i, j \le n, \text{ so } A+B \in W$
                    \item [iii]: $\alpha \in \mathbb{F}, A \in W: \text{ then } (\alpha A)_{ij} = \alpha a_{ij} = \alpha a_{ji} = (\alpha A)_{ji}, \forall 1 \le i, j \le n, \text{ so } \alpha A \in W$
                \end{itemize}
                \par Thus $W$ is a subspace of $V$ ($W \le V$).